\section*{Capítulo \ref{ch:b1:curves} --- Curvas Planas}

\begin{solution}{exo:b1:curves:1}
$M(-t) = (t^2, -t^3)$: reflexão no eixo $x$; estude $t \geq 0$. Tanto $x' = 2t$ quanto
$y' = 3t^2$ são $\geq 0$: o ramo move-se para a direita e para cima, de
$(0,0)$ ao infinito. Em $t = 0$ a velocidade se anula; as expansões
$x = t^2$, $y = t^3$ mostram que a tangente é o eixo $x$ ($y/x = t \to 0$), com
$y$ mudando de sinal enquanto $x \geq 0$: uma \emph{\omterm{ex:b1:curves:astroid}{cúspide}} apontando
para a esquerda. A curva é a parábola semicúbica $y^2 = x^3$.
\end{solution}

\begin{solution}{exo:b1:curves:2}
$M(t + 2\pi) = M(t) + (2\pi, 0)$: a curva se repete, transladada por uma circunferência de roda;
estude um período. $x'(t) = 1 - \cos t \geq
0$, $y'(t) = \sin t$: o arco sobe em $\intoo{0}{\pi}$, desce em
$\intoo{\pi}{2\pi}$, culminando em $(\pi, 2)$. Pontos singulares onde $x' = y' = 0$: $t \in 2\pi\Z$, no
solo. Perto de $t = 0$:
\[
x(t) = \frac{t^3}{6} + o(t^3),
\qquad
y(t) = \frac{t^2}{2} + o(t^2):
\]
$x$ muda de sinal, $y \geq 0$, e $\frac{x}{y^{3/2}}$ é limitado: a tangente é
\emph{vertical} (direção $(0,1)$), uma \omterm{ex:b1:curves:astroid}{cúspide} em que o ponto
acompanhado tem momentaneamente velocidade nula --- a assinatura
física do rolamento sem escorregamento.
\end{solution}

\begin{solution}{exo:b1:curves:3}
Em $t = \frac\pi4$: $M = \bigl(\frac{\sqrt2}{4},
\frac{\sqrt2}{4}\bigr)$, direção tangente $(-\cos t, \sin t) =
\frac{1}{\sqrt2}(-1, 1)$ (o \cref{ex:b1:curves:astroid}). Reta tangente:
$y - \frac{\sqrt2}{4} = -(x - \frac{\sqrt2}{4})$, isto é, $x +
y = \frac{\sqrt2}{2}$. Interseções com os eixos: $\bigl(\frac{\sqrt2}{2},
0\bigr)$ e $\bigl(0, \frac{\sqrt2}{2}\bigr)$; o
segmento entre eles tem \omterm{def:b1:curves:arclength}{comprimento} $\sqrt{\frac12 + \frac12} = 1$.

(Ponto geral $t$: a tangente em $(\cos^3 t, \sin^3 t)$ corta os eixos em $(\cos t, 0)$ e $(0, \sin t)$
--- verifique que a reta por esses pontos tem direção $(-\cos t, \sin t)$ e passa
por $M(t)$ --- e o segmento cortado tem \omterm{def:b1:curves:arclength}{comprimento} $\sqrt{\cos^2 t +
\sin^2 t} = 1$.)
\end{solution}

\begin{solution}{exo:b1:curves:4}
$r = \cos 2\theta$: período $\pi$ em $\theta$, simétrica nos dois eixos; $r$
anula-se em $\theta = \pm\frac\pi4$ (tangentes na origem ao longo das diagonais) e é
negativo para $\theta \in
\intoo{\frac\pi4}{\frac{3\pi}{4}}$, onde os pontos são marcados no raio oposto ---
produzindo quatro pétalas ao longo das direções dos eixos $\theta = 0, \frac\pi2, \pi, \frac{3\pi}{2}$, cada
uma de raio máximo $1$.

$r\cos\theta = 1$ é a equação $x = 1$: a \omterm{def:b1:curves:polar}{curva polar} $r =
\frac{1}{\cos\theta}$ é a reta vertical
$x = 1$ (percorrida uma vez para $\theta \in \intoo{-\frac\pi2}{\frac\pi2}$).
\end{solution}

\begin{solution}{exo:b1:curves:5}
$r(\theta) = 1 + 2\cos\theta$ anula-se para $\cos\theta = -\frac12$: $\theta = \pm\frac{2\pi}{3}$. Simetria no eixo $x$; estude $\theta \in \intcc{0}{\pi}$.
Para $\theta \in
\intco{0}{\frac{2\pi}{3}}$: $r > 0$, decrescente de $3$ a $0$: arco externo, entrando
na origem tangente ao raio $\theta =
\frac{2\pi}{3}$. Para $\theta \in \intoc{\frac{2\pi}{3}}{\pi}$: $r < 0$: os pontos
$r\vec u(\theta)$ ficam no raio oposto (ângulo $\theta - \pi \in \intoc{-\frac\pi3}{0}$), com distância $\abs r$
crescendo de $0$ a $1$: isto desenha um pequeno \emph{laço interno}
pela origem. Em $\theta = \pi$, $r = -1$ e o ponto é $-\vec u(\pi) = (1, 0)$: o laço fecha-se
no semieixo $x$ positivo. Esboço: uma grande curva externa em forma
de coração, com alcance máximo $3$ em $\theta = 0$, mais um laço dentro
dela, por $O$ e por $(1,0)$.
\end{solution}

\begin{solution}{exo:b1:curves:6}
Complete os quadrados:
\[
9(x^2 - 4x) + 25(y^2 - 2y) = 164
\iff 9(x-2)^2 + 25(y-1)^2 = 164 + 36 + 25 = 225 .
\]
Dividindo por $225$: $\frac{(x-2)^2}{25} + \frac{(y-1)^2}{9} = 1$: uma elipse de centro $(2, 1)$, semieixos
$a = 5$ (horizontal) e $b = 3$; $c = \sqrt{25 - 9} = 4$: focos $(2 \pm 4,\, 1) = (-2, 1)$ e $(6, 1)$;
excentricidade $e = \frac45$.
\end{solution}

\begin{solution}{exo:b1:curves:7}
Foco na origem, diretriz $D: x = d$ com $d = \frac pe > 0$. Para um ponto $M = (r\cos\theta, r\sin\theta)$ com
$r > 0$:
\[
MF = r,
\qquad
d(M, D) = \abs{d - r\cos\theta},
\]
e a condição da \omterm{def:b1:curves:conics}{cônica} $MF = e\,d(M, D)$, no regime em que $M$ está do lado do
foco em relação à diretriz ($r\cos\theta < d$), lê-se $r
= e(d - r\cos\theta)$, isto é,
\[
r(1 + e\cos\theta) = ed = p,
\qquad
r = \frac{p}{1 + e\cos\theta} .
\]
Para $e < 1$ o denominador nunca se anula: todos os $\theta$ são
permitidos (elipse). Para $e = 1$: $\theta \neq \pi$ (parábola, \omterm{def:b1:topology:open}{aberta} para o
lado oposto ao da diretriz). Para $e > 1$: é preciso $1 + e\cos\theta > 0$, isto é,
$\theta \in \intoo{-\theta_0}{\theta_0}$ com $\theta_0 =
\arccos\bigl(-\frac1e\bigr)$: um \emph{ramo} da hipérbole (o outro ramo
corresponde à escolha de sinal $r < 0$, ou ao segundo foco).
\end{solution}

\begin{solution}{exo:b1:curves:8}
\emph{Jardineiro:} amarre um barbante de \omterm{def:b1:curves:arclength}{comprimento} $2a$ a duas
estacas $F, F'$ e mantenha-o esticado com o ponto traçador $M$: a
restrição é exatamente $MF + MF' = 2a$, logo a curva traçada é a elipse.

\emph{Propriedade de reflexão:} seja $t \mapsto M(t)$ uma parametrização regular
e $u(t) = \frac{M(t) - F}{\norm{M(t) - F}}$, sendo $u'(t)$ o vetor unitário análogo em direção a $F'$.
Derivando $\norm{M - F} + \norm{M - F'} = 2a$, usando $\frac{\dd}{\dd
t}\norm{M - F} = \bigl\langle \frac{M - F}{\norm{M-F}},\,
M'\bigr\rangle$ (regra da cadeia em $\sqrt{\langle\cdot,\cdot\rangle}$):
\[
\bigl\langle u + u',\, M'\bigr\rangle = 0 .
\]
Logo a direção tangente $M'$ é \omterm{def:b1:euclid:orthogonal}{ortogonal} à direção bissetriz $u + u'$
dos dois raios focais: a tangente faz ângulos iguais com $MF$ e
$MF'$. (Um raio de luz que parte de um foco reflete-se na elipse para
o outro foco.)
\end{solution}

\begin{solution}{exo:b1:curves:9}
\begin{enumerate}
  \item Domínio: $\cos 2\theta \geq 0$: $\theta \in
        \intcc{-\frac\pi4}{\frac\pi4} \cup
        \intcc{\frac{3\pi}{4}}{\frac{5\pi}{4}}$. Simetrias: $\theta
        \mapsto -\theta$ (eixo $x$) e $\theta \mapsto \pi -
        \theta$
        (eixo $y$): estude $\intcc{0}{\frac\pi4}$; $r$ decresce de $1$ a $0$. Em
        $\theta = \pm\frac\pi4$: $r =
        0$, tangentes na origem ao longo das diagonais. A
        curva é o símbolo $\infty$: dois laços simétricos que se
        encontram em $O$, atingindo $(\pm 1, 0)$.
  \item Com $F, F' = (\pm c, 0)$, $c = \frac{1}{\sqrt 2}$, e $M
        = (r\cos\theta, r\sin\theta)$:
        \[
        MF^2\, MF'^2
        = \bigl((r\cos\theta - c)^2 + r^2\sin^2\theta\bigr)
        \bigl((r\cos\theta + c)^2 + r^2\sin^2\theta\bigr)
        = (r^2 + c^2)^2 - (2rc\cos\theta)^2 ,
        \]
        pela identidade $(A - B)(A + B) = A^2 - B^2$ com $A = r^2
        + c^2$, $B = 2rc\cos\theta$. Com $c^2 = \frac12$:
        \[
        MF^2 MF'^2 = \Bigl(r^2 + \frac12\Bigr)^2 - 2r^2\cos^2\theta
        = r^4 + r^2\bigl(1 - 2\cos^2\theta\bigr) + \frac14
        = r^4 - r^2\cos 2\theta + \frac14 .
        \]
        Sobre a lemniscata, $r^2 = \cos 2\theta$: os dois primeiros termos se
        cancelam, restando $MF^2MF'^2 = \frac14$, isto é, $MF \cdot MF' =
        \frac12$. Reciprocamente, o
        cálculo lido ao contrário mostra que a equação do lugar $MF\,MF' = \frac12$
        é $r^2 = \cos
        2\theta$ (para $r \neq 0$; e $O$ satisfaz as duas).
\end{enumerate}
\end{solution}

\begin{solution}{exo:b1:curves:10}
$r = 1 + \cos\theta$, $r' = -\sin\theta$:
\[
r^2 + r'^2 = 1 + 2\cos\theta + \cos^2\theta + \sin^2\theta
= 2 + 2\cos\theta = 4\cos^2\frac\theta2 ,
\]
logo $\sqrt{r^2 + r'^2} = 2\abs{\cos\frac\theta2}$ e, pela simetria no eixo $x$,
\[
L = \int_0^{2\pi} 2\Bigl|\cos\frac\theta2\Bigr|\,\dd\theta
= 2\int_0^{\pi} 2\cos\frac\theta2\,\dd\theta
= 2\Bigl[4\sin\frac\theta2\Bigr]_0^{\pi} = 8 .
\]
Mais um perímetro algébrico, sem $\pi$.
\end{solution}

\begin{solution}{exo:b1:curves:11}
O ponto $(a\cos t, b\sin t)$ satisfaz a equação: $\cos^2 t + \sin^2 t = 1$. O vetor normal da reta é
$\bigl(\frac{\cos t}a, \frac{\sin t}b\bigr)$, e o seu produto com a velocidade $(-a\sin t, b\cos t)$ é $-\sin t\cos t +
\sin t\cos t = 0$: a reta passa
pelo ponto com a direção tangente --- ela é a tangente. Para $(x_0, y_0)$
sobre a elipse, escreva $\cos t = \frac{x_0}a$, $\sin t = \frac{y_0}b$ e substitua:
\[
\frac{x\,x_0}{a^2} + \frac{y\,y_0}{b^2} = 1 ,
\]
obtida da equação da elipse ``desdobrando'' $x^2 \mapsto
x\,x_0$ e $y^2 \mapsto y\,y_0$.
\end{solution}

\begin{solution}{exo:b1:curves:12}
\begin{enumerate}
  \item $r' = k\eu^{k\theta}$, logo $\tan V = \frac{r}{r'} =
        \frac1k$: constante. A espiral corta todo raio a
        partir da origem sob o mesmo ângulo $V = \arctan\frac1k$.
  \item Em notação complexa a espiral é $\{\eu^{k\theta}
        \eu^{\iu\theta} : \theta \in \R\}$. Girar por $c$
        multiplica por $\eu^{\iu c}$:
        \[
        \eu^{k\theta}\eu^{\iu(\theta + c)}
        = \eu^{-kc}\,\eu^{k(\theta + c)}\eu^{\iu(\theta+c)} ,
        \]
        e, à medida que $\theta + c$ percorre $\R$, isto descreve $\eu^{-kc}$
        vezes a espiral: rotação $=$ homotetia. Nenhuma outra curva
        suave, além das retas e dos círculos, tem essa propriedade.
  \item $\sqrt{r^2 + r'^2} = \sqrt{1 + k^2}\,\eu^{k\theta}$, logo em $\intcc{A}{\theta_0}$ o \omterm{def:b1:curves:arclength}{comprimento} é $\frac{\sqrt{1+k^2}}{k}\bigl(\eu^{k\theta_0} -
        \eu^{kA}\bigr)$, e quando $A \to -\infty$:
        \[
        L = \frac{\sqrt{1 + k^2}}{k}\,\eu^{k\theta_0} < \infty :
        \]
        infinitas voltas em torno da origem, \omterm{def:b1:curves:arclength}{comprimento} total finito
        (as voltas encolhem geometricamente).
\end{enumerate}
\end{solution}

\begin{solution}{pb:b1:curves:1}
\textbf{1.} Rolar sem escorregar significa que o ponto de contato
percorreu uma distância igual ao arco de roda desenrolado: depois de
girar por $t$, o centro está em $(Rt, R)$. O ponto marcado fica no
aro, num ângulo $t$ \emph{atrás} da vertical descendente (a roda gira
no sentido horário enquanto avança):
\[
M(t) = \Omega(t) + R(-\sin t, -\cos t)
= \bigl(R(t - \sin t),\ R(1 - \cos t)\bigr),
\]
o que é correto em $t = 0$ ($M = (0,0)$) e em $t = \pi$ ($M =
(\pi R, 2R)$, o topo).

\textbf{2.} $f'(t) = R(1 - \cos t,\ \sin t)$ e
\[
\norm{f'}^2 = R^2\bigl((1-\cos t)^2 + \sin^2 t\bigr)
= 2R^2(1 - \cos t) = 4R^2\sin^2\frac t2 ,
\]
logo $\norm{f'} = 2R\abs{\sin\frac t2}$. Pontos singulares em $t \in
2\pi\Z$: as \omterm{ex:b1:curves:astroid}{cúspides} no solo
(o \cref{exo:b1:curves:2}); o topo do arco é $t = \pi$, onde a velocidade $2R$ é
máxima.

\textbf{3.} Fatorações de meio ângulo:
\[
f'(t) = 2R\sin\frac t2\,\Bigl(\sin\frac t2,\ \cos\frac t2\Bigr),
\quad
M - C = 2R\sin\frac t2\,\Bigl(-\cos\frac t2,\ \sin\frac t2\Bigr),
\]
\[
T - M = \bigl(R\sin t,\ R(1 + \cos t)\bigr)
= 2R\cos\frac t2\,\Bigl(\sin\frac t2,\ \cos\frac t2\Bigr).
\]
Para $0 < t < 2\pi$, $\sin\frac t2 \neq 0$: $M - C$ é \omterm{def:b1:euclid:orthogonal}{ortogonal} à direção tangente $\bigl(\sin\frac t2, \cos\frac
t2\bigr)$ (o
produto deles é $-\sin\frac t2\cos\frac t2 +
\sin\frac t2\cos\frac t2 = 0$), e $T - M$ é paralelo a ela. A corda até o
topo da roda \emph{é} a tangente; a corda até o ponto de contato é a
normal.

\textbf{4.} A cada instante a roda pivota em torno do seu ponto de
contato (esse ponto tem velocidade nula: rolamento sem
escorregamento), logo todo ponto rígido da roda move-se
\omterm{def:b1:euclid:orthogonal}{ortogonalmente} à reta que o liga a $C$, com velocidade escalar
(velocidade angular $1$) igual a essa distância. Verificação: $\norm{M - C} = 2R\abs{\sin\frac t2} =
\norm{f'(t)}$.

\textbf{5.} $2R\,y(t) = 2R\cdot 2R\sin^2\frac t2 =
4R^2\sin^2\frac t2 = \norm{f'(t)}^2$. A velocidade escalar à altura $y$ é $\sqrt{2R\,y}$ ---
formalmente a lei $v = \sqrt{2gh}$ da queda livre, com o solo fazendo o papel de
teto; a Parte IV transforma esta observação em maquinaria de relógio.

\textbf{6.} Pela \cref{def:b1:curves:arclength} e pela questão 2 ($\sin\frac t2 \geq 0$ em $\intcc{0}{2\pi}$):
\[
L = \int_0^{2\pi} 2R\sin\frac t2\,\dd t
= 2R\Bigl[-2\cos\frac t2\Bigr]_0^{2\pi} = 2R(2 + 2) = 8R .
\]
O teorema de Wren: exatamente quatro diâmetros.

\textbf{7.} $s(t) = \int_0^t 2R\sin\frac u2\,\dd u =
4R\bigl(1 - \cos\frac t2\bigr)$; $s(2\pi) = 4R(1+1) = 8R$, coerente.

\textbf{8.} $\sigma(t) = \abs{s(t) - s(\pi)} = \abs{4R(1 -
\cos\frac t2) - 4R} = 4R\abs{\cos\frac t2}$, logo $\sigma^2 =
16R^2\cos^2\frac t2$; e $2R - y = R(1 + \cos t) =
2R\cos^2\frac t2$, donde $8R(2R - y) = 16R^2\cos^2\frac t2 =
\sigma^2$.

\textbf{9.} Na \omterm{ex:b1:curves:astroid}{cúspide} $t = 0$ (ou $2\pi$): $\sigma = 4R$, metade de $8R$: o
ápice \omterm{def:b1:arith:divides}{divide} o arco ao meio. Nos dois cálculos a velocidade escalar é
$\abs{\text{polinômio trigonométrico em } t/2}$, cuja primitiva é de novo trigonométrica: o \omterm{def:b1:curves:arclength}{comprimento} é uma
diferença de \emph{valores} de cossenos --- números racionais vezes
$R$ --- sem nenhum arco de círculo a medir, logo sem $\pi$. O próprio
círculo tem velocidade constante, de modo que a sua \omterm{thm:b1:integration:def}{integral} de
\omterm{def:b1:curves:arclength}{comprimento} produz o \omterm{def:b1:curves:arclength}{comprimento} total do \omterm{prop:b1:reals:intervals}{intervalo}, $2\pi$; a
velocidade da cicloide anula-se nas extremidades e integra-se
algebricamente.

\textbf{10.} O arco é percorrido com $x$ crescendo de $0$ a
$2\pi R$ ($x' = R(1 - \cos t) \geq 0$, anulando-se apenas em pontos isolados). Substituir
$x = x(t)$ na \omterm{thm:b1:integration:def}{integral} de área $\int_0^{2\pi R} y\,\dd x$ (o \cref{thm:b1:integration:parts}) dá $A = \int_0^{2\pi} y(t)\,x'(t)\,\dd t$.

\textbf{11.}
\[
A = \int_0^{2\pi} R(1 - \cos t)\cdot R(1 - \cos t)\,\dd t
= R^2\int_0^{2\pi}\bigl(1 - 2\cos t + \cos^2 t\bigr)\dd t
= R^2\Bigl(2\pi - 0 + \pi\Bigr) = 3\pi R^2 ,
\]
usando $\int_0^{2\pi}\cos^2 = \pi$. Exatamente três áreas de roda: a balança de Galileu
dizia ``cerca de $3$''; o cálculo de Roberval diz ``exatamente''.

\textbf{12.} Com $x = \cos^3 t$, $y = \sin^3 t$: $y\,x' =
-3\sin^4 t\cos^2 t$. Linearize:
\[
\sin^4 t\cos^2 t = \frac{1 - \cos 2t}{2}\cdot\frac{\sin^2
2t}{4} = \frac{\sin^2 2t}{8} - \frac{\sin^2 2t\cos 2t}{8} ,
\]
e em $\intcc{0}{2\pi}$: $\int\sin^2 2t = \pi$, $\int \sin^2
2t\cos 2t = \bigl[\frac{\sin^3 2t}{6}\bigr] = 0$. Logo $\oint
y\,\dd x = -3\cdot\frac\pi8$, e a área encerrada é $\frac{3\pi}8$ (o
sinal registra a orientação anti-horária).

\textbf{13.} Para $(\cos t, \sin t)$ com $t$ indo de $\pi$ a $0$, $x$ cresce de
$-1$ a $1$ e
\[
\int_\pi^0 \sin t\cdot(-\sin t)\,\dd t
= \int_0^\pi \sin^2 t\,\dd t = \frac\pi2 :
\]
a fórmula devolve a área do semidisco superior, como deve.

\textbf{14.} $x' = R(1 + \cos t) = 2R\cos^2\frac t2$, $y' =
R\sin t = 2R\sin\frac t2\cos\frac t2$, logo $\norm{f'} =
2R\cos\frac t2$ (não negativa em $\intcc{-\pi}{\pi}$). Arco a
partir do fundo: $s(t) = \int_0^t 2R\cos\frac u2\,\dd u =
4R\sin\frac t2$. Então
\[
y = R(1 - \cos t) = 2R\sin^2\frac t2
= 2R\Bigl(\frac{s}{4R}\Bigr)^{\!2} = \frac{s^2}{8R} .
\]

\textbf{15.} Conservação da energia com $v = \frac{\dd s}{\dd\tau}$ e $y = \frac{s^2}{8R}$, $y_0 = \frac{s_0^2}{8R}$:
\[
\Bigl(\frac{\dd s}{\dd\tau}\Bigr)^{\!2} = 2g(y_0 - y)
= \frac{2g}{8R}\bigl(s_0^2 - s^2\bigr)
= \frac{g}{4R}\bigl(s_0^2 - s^2\bigr).
\]

\textbf{16.} Derivando em $\tau$: $2s's'' =
-\frac{g}{4R}\,2ss'$, logo, onde quer que $s' \neq 0$ (e
portanto em toda parte, por \omterm{def:b1:continuity:continuous}{continuidade}), $s'' = -\frac{g}{4R}s$: o oscilador
harmônico. Pelo \cref{thm:b1:diffeq:homogeneous2}, $s(\tau) =
A\cos\omega\tau + B\sin\omega\tau$ com $\omega =
\sqrt{g/(4R)}$; as condições iniciais $s(0) = s_0$,
$s'(0) = 0$ dão $s(\tau) = s_0\cos\omega\tau$.

\textbf{17.} A conta atinge o fundo quando $s = 0$, isto é, em $\omega\tau = \frac\pi2$:
\[
T_{\downarrow} = \frac{\pi}{2\omega}
= \frac\pi2\sqrt{\frac{4R}{g}} = \pi\sqrt{\frac Rg}\,,
\]
independente de $s_0$: soltas de qualquer lugar da tigela, todas as
contas chegam no mesmo instante --- a tautócrona.

\textbf{18.} Substituindo $s = s_0\cos\omega\tau$: o lado esquerdo é $s_0^2\omega^2\sin^2\omega\tau$ e o lado
direito é $\frac{g}{4R}s_0^2(1 - \cos^2\omega\tau) =
s_0^2\omega^2\sin^2\omega\tau$: igualdade exata, logo nenhuma solução espúria foi
introduzida. Para um arco circular, $y$ não é proporcional a $s^2$
($y = \ell(1 - \cos\frac s\ell) =
\frac{s^2}{2\ell} - \frac{s^4}{24\ell^3} + \dots$): o termo restaurador é apenas \emph{aproximadamente} \omterm{def:b1:linmaps:def}{linear},
de modo que o período de um pêndulo circular varia com a amplitude,
ao passo que o da cicloide é rigorosamente constante.

\textbf{19.} $\pi\sqrt{R/g} = 1$ dá $R = \frac{g}{\pi^2}
\approx \frac{9.81}{9.87} \approx 0.994$ m --- quase exatamente um metro. A
oscilação completa leva $\frac{2\pi}\omega =
2\pi\sqrt{4R/g}$, que é precisamente a fórmula de
pequenas oscilações $2\pi\sqrt{\ell/g}$ para um pêndulo de \omterm{def:b1:curves:arclength}{comprimento} $\ell = 4R$:
Huygens suspendeu o seu pêndulo entre faces cicloidais exatamente
dessa proporção.

\textbf{20.} Em $t = \frac\pi2$, $R = 1$: $M = \bigl(\frac\pi2 -
1,\ 1\bigr)$, $T = \bigl(\frac\pi2,\ 2\bigr)$, logo $T - M = (1,
1)$. E $f'(\frac\pi2) = (1 - 0,\ 1) = (1, 1)$: a
corda até o topo é a velocidade, como prometido.

\textbf{21.} Para a trocoide, $x'(t) = R - d\cos t$ e $y'(t)
= d\sin t$. Se $d < R$: $x' \geq R - d > 0$, o
ponto sempre avança; a velocidade nunca se anula (a sua primeira
componente é positiva): nenhum ponto singular, uma onda regular. Se
$d > R$: $x'(0)
= R - d < 0 < R + d = x'(\pi)$, logo o ponto move-se para trás perto dos instantes de
contato e para a frente nos demais: a curva cruza-se a si mesma em
laços. Um ponto no friso de uma roda ferroviária (abaixo do boleto do
trilho, $d > R$) viaja para trás a cada volta.

\textbf{22.} A corda da \omterm{ex:b1:curves:astroid}{cúspide} $(\pi R, 2R)$ até a origem tem \omterm{def:b1:curves:arclength}{comprimento}
$L = R\sqrt{\pi^2 + 4}$ e ângulo de inclinação $\alpha$ com $\sin\alpha = \frac{2R}{L}$. Deslizando do
repouso com aceleração constante $g\sin\alpha$: $L = \frac12 g\sin\alpha\,T^2$, logo
\[
T = \sqrt{\frac{2L}{g\sin\alpha}} = \sqrt{\frac{2L^2}{2Rg}}
= \frac{L}{\sqrt{Rg}} = \sqrt{\pi^2 + 4}\,\sqrt{\frac Rg}
\approx 3.72\sqrt{\frac Rg}\,,
\]
contra $\pi\sqrt{R/g} \approx 3.14\sqrt{R/g}$ para a cicloide: o caminho curvo é mais rápido. Ele é,
de fato, o mais rápido possível --- a braquistócrona --- um teorema
do cálculo das variações, além deste volume.

\textbf{23.} $\dfrac{\dd y}{\dd x} = \dfrac{y'}{x'} =
\dfrac{\sin t}{1 - \cos t} = \cot\dfrac t2$ (fórmulas de meio ângulo). Então
\[
\frac{\dd^2y}{\dd x^2}
= \frac{\dd}{\dd t}\Bigl(\cot\frac t2\Bigr)\cdot\frac{1}{x'(t)}
= -\frac{1}{2\sin^2\frac t2}\cdot\frac{1}{2R\sin^2\frac t2}
= -\frac{1}{4R\sin^4\frac t2} < 0 :
\]
côncavo ao longo de todo o arco; quando $t \to 0^+$ ou $2\pi^-$, o
coeficiente angular $\cot\frac t2 \to \pm\infty$: tangentes verticais nas \omterm{ex:b1:curves:astroid}{cúspides}.

\textbf{24.} A direção da corda é $T - M \parallel
\bigl(\sin\frac t2, \cos\frac t2\bigr)$, que tende a $(0, 1)$ quando
$t \to 0^{+}$: a tangente na \omterm{ex:b1:curves:astroid}{cúspide} é vertical --- recuperada por pura
geometria, sem nenhuma expansão de Taylor.

\textbf{25.} (i) O teorema de Wren, $L = 8R$; a área de Roberval,
$A =
3\pi R^2$ (a razão $3$ conjecturada por Galileu); a tautócrona de
Huygens, $T_\downarrow = \pi\sqrt{R/g}$. (ii) Todo cálculo apoiou-se nas fatorações de meio
ângulo $f' =
2R\sin\frac t2(\sin\frac t2, \cos\frac t2)$ e no seu análogo para a tigela --- uma única identidade
alimentando tangente, \omterm{def:b1:curves:arclength}{comprimento} e relógio. (iii) A relação
intrínseca $y = s^2/8R$ converteu a equação de energia em $s'' = -\frac{g}{4R}s$, uma equação
\omterm{def:b1:linmaps:def}{linear} de coeficientes constantes cujas soluções são exatamente
isócronas. (iv) A Parte I usou o cálculo diferencial do \cref{ch:b1:derivative}, as
Partes II--III a \omterm{thm:b1:integration:def}{integral} do \cref{ch:b1:integration}, e a Parte IV as equações
diferenciais do \cref{ch:b1:diffeq} --- a cicloide é o currículo deste livro
enrolado numa só curva.
\end{solution}
